% arara: xelatex: { options: ["--synctex=1", "-interaction=nonstopmode", "-output-directory=../build"] }
% arara: move: { files: ["../build/summary_3page.pdf"], target: "../pdf/" }

\documentclass[aspectratio=169,10pt]{beamer}
\usepackage{graphicx}
\usepackage{amssymb,amsfonts,amsmath}
\usepackage{xeCJK}
\usepackage{tcolorbox}
\usepackage{tikz}
\usepackage{booktabs}
\usetikzlibrary{positioning,shapes.geometric,arrows.meta,calc}
\setCJKmainfont{STSong}
\setCJKsansfont{STHeiti}
\setCJKmonofont{STFangsong}
\usefonttheme[onlymath]{serif}

% 自定义颜色
\definecolor{highlightyellow}{RGB}{255, 242, 0}
\definecolor{dotred}{RGB}{200, 30, 30}
\definecolor{darkgreen}{RGB}{0, 128, 0}
\definecolor{layer1}{RGB}{227,242,253}
\definecolor{layer2}{RGB}{187,222,251}

\title{突破 Coding Agent 能力天花板\\{\large 基于 Context Graph 的记忆架构演进}}
\author{WU Zihan}
\date{\today}

\begin{document}

\begin{frame}
  \maketitle
\end{frame}

% Page 1
\begin{frame}
  \frametitle{1. 业界现状：从"长窗口"到"结构化记忆"的演进}

  \begin{columns}[T]
    \begin{column}{0.48\textwidth}
      \begin{block}{主流模型记忆机制 (State of the Art)}
        \footnotesize
        \begin{itemize}
            \item \textbf{超长上下文 (Brute Force)}:
            \newline Gemini (1M+), Claude (200K)。
            \newline \textit{局限：被动且昂贵，无法跨会话累积。}
            \item \textbf{向量 RAG (Vector Database)}:
            \newline ChatGPT, Auto-GPT。基于语义相似度检索。
            \newline \textit{局限：\textcolor{dotred}{丢失结构与关系}}（知道"A"和"B"相似，但不知"A导致B"）。
            \item \textbf{混合/隐式图谱}:
            \newline Gemini "Personal Intelligence" (隐式关联用户数据)。
        \end{itemize}
      \end{block}

      \vspace{0.2cm}

      \begin{alertblock}{关键痛点：被动检索 vs 主动总结}
        \footnotesize
        传统 RAG \textbf{被动挖掘}已有知识，无法\textbf{主动抽象}经验规律。
        \newline 资深工程师的价值 = "踩过这个坑" + "总结了通用解法"。
        \newline \textcolor{dotred}{LLM 需要主动记忆机制，而非每次重新检索}。
      \end{alertblock}
    \end{column}

    \begin{column}{0.48\textwidth}
      \begin{block}{技术新趋势：向有向图 (Directed Graph) 转化}
        \footnotesize
        \textbf{趋势：Unstructured Vectors $\to$ Structured Graphs}
        \begin{itemize}
            \item \textbf{Microsoft GraphRAG}: 利用社区检测（Community Detection）解决全局性问题。
            \item \textbf{Zep (Temporal KG)}: 引入"时序知识图谱"，记忆事实的演变。
            \item \textbf{LangGraph}: 用图（StateGraph）规划 Agent 的思考流/控制流。
        \end{itemize}

        \vspace{0.2cm}

        \textbf{核心判断}：
        Coding 场景必须结合 \textbf{AST（依赖图）} + \textbf{执行流（因果图）}，结构化记忆是唯一解。
      \end{block}
    \end{column}
  \end{columns}
\end{frame}

% Page 2
\begin{frame}
  \frametitle{2. 核心方案：Context Graph 可生长记忆架构}

  \begin{tcolorbox}[colback=layer1, title=\textbf{方案本质：融合"数据图谱"与"控制流"的专用记忆体}]
    \small
    借鉴 \textbf{Zep (记忆存储)} 与 \textbf{LangGraph (流程控制)} 理念，构建专用于 Coding 的 \textbf{有向图谱}。
  \end{tcolorbox}

  \vspace{0.2cm}

  \begin{columns}[T]
    \begin{column}{0.55\textwidth}
      \begin{block}{架构设计：显式化"经验路径"}
        \footnotesize
        \textbf{1. 节点 (Nodes) - 沉淀痕迹}
        \begin{itemize}
            \item \texttt{Action/State}: 执行与反馈（类似 LangGraph 状态）。
            \item \texttt{Insight}: \textbf{经验胶囊}（\textcolor{darkgreen}{主动抽象}的通用规则，如"遇到X错误 $\to$ 修改Y配置"）。
        \end{itemize}

        \vspace{0.1cm}

        \textbf{2. 有向边 (Directed Edges) - 锁定逻辑}
        \begin{itemize}
            \item \texttt{NEXT} (时序): 记录操作流，复用成功路径。
            \item \texttt{CAUSES} (因果): 根因分析，\textbf{解决 Vector RAG 无法推理因果的缺陷}。
            \item \texttt{DEPENDS\_ON}: 代码静态依赖（AST）。
        \end{itemize}
      \end{block}
    \end{column}

    \begin{column}{0.42\textwidth}
      \begin{block}{差异化优势 (The Graph Advantage)}
        \footnotesize
        \begin{itemize}
            \item \textbf{主动学习 vs 被动检索}：
            \newline \textcolor{darkgreen}{主动总结}成功/失败模式，形成可复用规则。被动RAG只能"找相似案例"。

            \item \textbf{知识迁移能力 (Zero-Shot)}：
            \newline 图谱足够丰富时，能通过\textbf{路径推理}解决\textcolor{darkgreen}{未见过的问题}（如组合已知模式）。

            \item \textbf{数据飞轮}：
            \newline 任务越多 $\to$ 经验越丰富 $\to$ 解决新问题越快。
        \end{itemize}
      \end{block}
    \end{column}
  \end{columns}
\end{frame}
% Page 3
\begin{frame}
  \frametitle{3. 实施路线与业务价值}

  \begin{columns}[T]
    \begin{column}{0.58\textwidth}
      \begin{block}{落地路线图 (Roadmap)}
        \small
        \begin{enumerate}
            \item \textbf{Phase 1: MVP (2-3个月) —— 主动学习闭环}
            \begin{itemize}
                \item 场景：SWE-bench Lite。
                \item 目标：Agent \textcolor{darkgreen}{主动总结}失败模式，抽象为Insight节点，不再重复犯错。
            \end{itemize}

            \item \textbf{Phase 2: 企业级适配 (4-6个月) —— 存量攻坚}
            \begin{itemize}
                \item 场景：C/C++ 大型遗留系统。
                \item 技术：引入 AST 代码依赖分析，构建项目级图谱。
            \end{itemize}

            \item \textbf{Phase 3: 人机协同 (长期) —— 知识迁移}
            \begin{itemize}
                \item \textcolor{darkgreen}{跨项目迁移}：成熟图谱迁移到新项目，\textbf{零样本解决未见过的问题}。
                \item 吸纳人类工程师调试路径，持续丰富图谱。
            \end{itemize}
        \end{enumerate}
      \end{block}
    \end{column}

    \begin{column}{0.38\textwidth}
      \begin{alertblock}{业务价值 (Value)}
        \small
        \begin{itemize}
            \item \textbf{降本}：
            精准剪枝无效路径，\textbf{Token 节约 30\%}
            \item \textbf{增效}：
            阻断错误尝试，复杂 Bug \textbf{解决率 +20-30\%}
            \item \textbf{资产化}：
            消耗的算力转化为永久的\textbf{企业私有资产}（可复用的项目活文档）。
        \end{itemize}
      \end{alertblock}
    \end{column}
  \end{columns}

  \vspace{0.3cm}
  \centering
  \textbf{\small 我们构建的不是工具，而是具备"经验记忆"的 Autonomous Engineer。}
\end{frame}

% Page 4
\begin{frame}
  \frametitle{4. 业界实践与趋势：从向量到图谱的演进}

  \begin{columns}[T]
    \begin{column}{0.48\textwidth}
      \begin{block}{主流产品记忆机制现状}
        \footnotesize
        \begin{itemize}
          \item \textbf{ChatGPT}: 跨会话向量检索（2025）。
          \textcolor{dotred}{局限：丢失关系链。}

          \item \textbf{Claude}: 100K 超长上下文 + 知识库。

          \item \textbf{Gemini}: 1M tokens + Personal Intelligence（整合 Gmail/Photos）。

          \item \textbf{开源 (LangChain/Auto-GPT)}: 向量数据库。
        \end{itemize}
      \end{block}

      \vspace{0.1cm}

      \begin{alertblock}{核心痛点}
        \footnotesize
        向量 RAG 找"相似内容"，但\textbf{不知因果关系}，无法多跳推理。
      \end{alertblock}
    \end{column}

    \begin{column}{0.48\textwidth}
      \begin{block}{Context Graph 新趋势}
        \footnotesize
        \textbf{业界转向 Graph-Enhanced Memory：}

        \begin{itemize}
          \item \textbf{Zep/Graphiti}: 时序知识图谱，记录事实演变。准确率 94.8\%，延迟 -90\%。

          \item \textbf{LangGraph}: 控制流图（StateGraph）。

          \item \textbf{GraphRAG}: 社区检测 + 层级摘要。

          \item \textbf{Mem0}: 向量+图混合，速度 +91\%。
        \end{itemize}
      \end{block}

      \vspace{0.1cm}

      \begin{tcolorbox}[colback=layer1, title=\textbf{趋势判断}]
        \footnotesize
        \textbf{被动检索} $\to$ \textbf{主动学习}
        \begin{itemize}
          \item \textcolor{darkgreen}{主动总结}经验 vs 被动挖掘知识
          \item 确定性推理 vs 概率匹配
          \item 知识迁移：解决未见过的问题
        \end{itemize}
      \end{tcolorbox}
    \end{column}
  \end{columns}
\end{frame}

\end{document}